\subsection{数据加载：把序列切成 batch}

语言模型的数据本质上是一串整数 token。最常见的做法是把它序列化到 numpy 数组，再用 `memmap` 懒加载。

\begin{lstlisting}
start_indices = torch.randint(len(data) - sequence_length, (batch_size,))
x = torch.tensor([data[start:start + sequence_length] for start in start_indices])
\end{lstlisting}

这段逻辑的语义很直接：
\begin{itemize}
    \item 随机采样若干起点；
    \item 从每个起点切出一段固定长度的 token 序列；
    \item 堆成一个 batch。
\end{itemize}

\begin{knowledgebox}{为什么要用 memmap}
真实数据集可能非常大，不能一次性全读进内存。`np.memmap` 允许只把当前访问到的部分映射到内存中，因此更适合大规模语言建模。
\end{knowledgebox}

